%--- E2: after Lemma 2.5 and the paragraph on G_u, J_u (Section 2.2) ---------
\begin{explained}{annihilators, ideals, (F1)--(F3) and the splitting of the gcd}
\textbf{Step 1: annihilating a function.} Take $u(x)=e^{x}$. Since $\partial_xe^x=e^x$, $(\partial_x-1)e^x=0$: the operator $\partial_x-1$ \emph{annihilates} $e^x$. It is not the only one: $(\partial_x^2-1)e^x=0$, and $\partial_x^2-1=(\partial_x-1)(\partial_x+1)$ is annihilating precisely because it contains the factor $\partial_x-1$. More generally, if $Q(\partial_x)e^x=0$ then every polynomial multiple of $Q$ annihilates $e^x$, for instance $(\partial_x+3)(\partial_x-1)e^x=0$. This is the basic reason annihilators are ideals.

\medskip\textbf{Step 2: the annihilator is the whole ideal $(\partial_x-1)$.} For $R=\C[\partial_x]$, any polynomial can be divided by $z-1$: $Q(z)=(z-1)S(z)+Q(1)$. Replacing $z$ by $\partial_x$ and applying to $e^x$ gives $Q(\partial_x)e^x=Q(1)e^x$, hence
\[
Q(\partial_x)e^x=0\iff Q(1)=0\iff(z-1)\mid Q(z).
\]
So $\Ann(e^x)=(\partial_x-1)=\{S(\partial_x)(\partial_x-1):S\in\C[\partial_x]\}$, and every annihilating polynomial is a multiple of $\partial_x-1$.

\medskip\textbf{Step 3: why an ideal.} An ideal is a subset closed under addition (and additive inverses) and under multiplication by any element of the ring. For annihilators: if $Qu=0$ and $Pu=0$ then $(Q+P)u=Qu+Pu=0$; if $Qu=0$ and $S\in\ring$ then $(SQ)u=S(Qu)=S(0)=0$. Intuition: once an operator kills a function, composing it with another constant-coefficient operator cannot bring the function back to life.

\medskip\textbf{Step 4: one variable versus two variables.} In $\C[\xi]$ every ideal is generated by a single polynomial, $I=(g)$ -- like the multiples of $6$ in $\Z$. In $\C[\tau,\xi]$ several generators may be needed. Consider $I=(\tau,\xi-1)=\{\tau A+(\xi-1)B\}$: neither generator is a multiple of the other and the ideal is not generated by one polynomial. This matters because the annihilator of a function of two variables can have several independent generators: for $u(t,x)=e^x$ we have $\partial_tu=0$ and $(\partial_x-1)u=0$, and
\[
\Ann(u)=(\partial_t,\ \partial_x-1).
\]
Both equations are needed to describe it in the two-variable ring.

\medskip\textbf{Step 5: the three facts.}
\begin{itemize}[nosep]
\item \textbf{(F1) Unique factorization.} Every nonzero polynomial factors into irreducibles, uniquely up to constants and order; e.g.\ $\tau^2-\xi^2=(\tau-\xi)(\tau+\xi)$. This allows one to define the gcd of an ideal: the greatest common polynomial factor of all its elements.
\item \textbf{(F2) B\'ezout.} Two polynomials in $(\tau,\xi)$ with no common nonconstant factor have finitely many common zeros in $\C^2$. Example: $\tau$ and $\xi-1$ have no common factor and their only common zero is $(0,1)$. Geometrically, two algebraic curves without a shared component meet in finitely many points.
\item \textbf{(F3) Nullstellensatz.} If $q$ vanishes on $\Var(I)$ then $q^k\in I$ for some $k\ge1$: $q|_{\Var(I)}=0\Rightarrow q^k\in I$. It connects the geometry of a zero set to the algebra of the ideal. The power matters: $q$ itself need not belong to $I$ (e.g.\ $I=(\tau^2)$, $q=\tau$).
\end{itemize}

\medskip\textbf{Step 6: the splitting lemma.} Every nonzero ideal $I\subset\ring$ can be written $I=gJ$, with $g=\Gcd(I)$ and $\Var(J)$ finite. Here $g$ is the polynomial factor shared by every element of $I$; $J$ is the ideal of the quotients $f/g$; because $g$ is the gcd, the elements of $J$ have no common nonconstant factor, and by (F2) their common zero set is finite. So the ideal splits into a common polynomial factor and a residual ideal with finite zero set.
\begin{itemize}[nosep]
\item \emph{Example A (shared factor).} $I=(\tau(\xi-1),\tau(\xi+1))$. Both generators contain $\tau$, so $g=\tau$ and $J=(\xi-1,\xi+1)$. Since $(\xi+1)-(\xi-1)=2$, $J=(1)=\ring$, whose zero set is empty (hence finite). Thus $I=(\tau)$.
\item \emph{Example B (no shared factor).} $I=(\tau,\xi-1)$: $g=1$, $J=I$, $\Var(J)=\{(0,1)\}$, finite.
\end{itemize}

\medskip\textbf{Step 7: applied to $u$.} Define $G_u=\Gcd\Ann(u)$ and call the remaining ideal $J_u$, so $\Ann(u)=G_uJ_u$. For $u=e^x$ in $\C[\partial_t,\partial_x]$ we found $\Ann(e^x)=(\partial_t,\partial_x-1)$; the generators have no common nonconstant factor, so
\[
G_u=1,\qquad J_u=(\partial_t,\partial_x-1),\qquad \Var(J_u)=\{(0,1)\},
\]
the single point corresponding to the exponential $e^{0\cdot t+1\cdot x}=e^x$. A point of the characteristic picture corresponds to a possible exponential mode. \emph{Qualification:} the reading of $G_u$ as the ``genuinely infinite-dimensional part'' and of $J_u$ as the ``finitely many exponentials present'' is an intuition for the later theory, not a universal identification for arbitrary functions without further hypotheses.
\end{explained}

